% paper/sections/04c_dccd_derivation.tex
%
% §4.4  DCCD の導出形 + 実装デフォルト：1 次風上 → 修正方程式 → spectral filter
%        本稿での DCCD は CLS 移流の実装デフォルト（§7）；
%        bulk 移流 default は UCCD6（L1/L2/L3）；DCCD は CLS 移流 + reinit 圧縮項に適用．
%
% ═══════════════════════════════════════════════════════════════════════════════
\subsection{\DCCD{} フィルタ：導出形と実装デフォルト}
\label{sec:dccd_derivation}
\phantomsection\label{sec:dissipative_ccd}% backward-compat: 旧 04d §4.4 overall label
\phantomsection\label{sec:dccd_motivation}% backward-compat: 旧 04d § motivation
\phantomsection\label{sec:dccd_bc}% backward-compat: 旧 04d § BC (現在は §\ref{sec:ccd_bc} 参照)
\phantomsection\label{sec:dccd_conservation}% backward-compat: 旧 04d § conservation (本節最終項)
% ═══════════════════════════════════════════════════════════════════════════════

§~\ref{sec:ccd_bc}--§~\ref{sec:ccd_matrix} で確立した
基底 CCD のブロック Thomas 行列に対する
\textbf{散逸 post-filter} として DCCD を導出する．
本節の主役は基底 CCD の解 $(f, f', f'')$ ではなく，
それに作用する散逸演算子そのものである．

\medskip\noindent\textbf{なぜ基底 CCD に散逸が必要か．}
基底 CCD（§~\ref{sec:ccd_def}）は中心差分系であり，
最高波数 $\xi=k h\to\pi$ におけるシンボル
$\omega_1(\xi),\omega_2(\xi)$ の挙動から，
\textbf{Nyquist 近傍のチェッカーボード（$2h$ 波長振動）が散逸ゼロで温存される}．
これは平滑場の純伝搬では問題ないが，
実用シミュレーションでは離散化誤差・非線形相互作用・初期データ不規則性が
このモードを励起し，スプリアス振動として顕在化する．

\medskip\noindent\textbf{導出形と実装デフォルトの 2 段構成．}
本節では DCCD を以下の 2 段で示す：
\begin{itemize}[leftmargin=*, itemsep=2pt, topsep=2pt]
  \item \textbf{導出形}（§\ref{sec:dccd_embedding}）：
        中心 Chu--Fan CCD の 1 次導関数 $D_1^{\CCD}$ は純虚数 Fourier シンボルで散逸ゼロのため，
        \textbf{実数 Hermite 正定値}な偶数次演算子 $-D_2^{\CCD}$ を散逸チャンネルとして並列付加する筋道．
        DCCD の散逸チャンネルが「アドホックな数値粘性」ではなく
        中心系の対称性と CCD が既に解いている量の有効再利用であることを明らかにする．
  \item \textbf{実装デフォルト}（§\ref{sec:dccd_spectral_filter}--§\ref{sec:dccd_adaptive}）：
        伝達関数 $H(\xi)=1-4\varepsilon_d\sin^2(\xi/2)$ による
        スペクトル post-filter 形式．
        導出形 PDE はその低波数極限に対応する．
\end{itemize}
DCCD は本章の派生スキームのうち\textbf{基底 CCD への外殻 post-filter} の
形式を取る唯一のものであり，
散逸を演算子内部に埋め込む UCCD6（§~\ref{sec:uccd6_def}）と
本節の post-filter 形式は対照をなす．
具体的な用途章割当は §~\ref{sec:dccd_eps_design} の
\textbf{適用範囲の制限}で再述するに留める．

% -------------------------------------------------------------------------------
\subsubsection{Chu--Fan CCD への散逸チャンネルの埋込み}
\label{sec:dccd_embedding}
\phantomsection\label{sec:dccd_modified_eq}% backward-compat: 旧 §4.7.1 削除に伴う参照保護
% -------------------------------------------------------------------------------

中心 Chu--Fan CCD（§\ref{sec:ccd_def}）の 1 次導関数 $D_1^{\CCD}$ は
Fourier シンボル $i\omega_1(\xi)/h$ で\textbf{純虚数}（中心ステンシルの対称性に由来）であり，
固有値の実部がゼロで散逸を持たない．
散逸を加えるには\textbf{実数の Hermite 演算子}を主項に並列付加する必要がある．

基底 CCD が既に解いている 2 次導関数 $D_2^{\CCD}$ は
Fourier シンボル $-\omega_2(\xi)^2/h^2$ で\textbf{実かつ負定値}
（§\ref{sec:ccd_omega_weights}）であり，
$-D_2^{\CCD}$ の符号反転で正値の散逸チャンネルが得られる．
これを $\varepsilon\,|c|\,\Delta x$ のスケーリングで主項に並列して付加：
\begin{equation}
  \boxed{\;
  \partial_t u_i + c\,(D_1^{\CCD}u)_i
  \;=\; \varepsilon\,|c|\,\Delta x\,(D_2^{\CCD}u)_i,
  \qquad \varepsilon > 0.
  \;}
  \label{eq:dccd_pde_form}
\end{equation}
スケーリング $|c|\Delta x$ は次元解析（散逸係数の物理次元 $[L^2/T]$）から定まり，
$\varepsilon$ は無次元の振幅パラメタである．
CCD 自身の $D_2^{\CCD}$ を用いることが肝要で，
低次有限差分で代用するとステンシル整合性が崩れ，
修正波数の不整合が散逸チャンネルの精度を $\Ord{h^2}$ まで劣化させる．
本式は中心系の対称性のみから自閉的に導出され，
1 次風上スキーム等の外部参照を要しない．
（風上ステンシル + ハイパー粘性として散逸を演算子内部に組み込む別経路は
UCCD6（§\ref{sec:uccd6_def}）で展開する．）

% -------------------------------------------------------------------------------
\subsubsection{実装デフォルト：スペクトルフィルタとの関係}
\label{sec:dccd_spectral_filter}
\phantomsection\label{sec:dccd_filter_theory}% backward-compat: 旧 04d § filter theory
% -------------------------------------------------------------------------------

実装デフォルト DCCD は定数 $\varepsilon$ を
スペクトル伝達関数：
\begin{equation}
  H(\xi) = 1 - 4\varepsilon_d\sin^2(\xi/2),
  \qquad \xi=k\Delta x,
  \label{eq:dccd_filter_transfer}
\end{equation}
% backward-compat phantom labels removed (CHK-205 m2): use eq:dccd_filter_transfer
に置き換える．
低波数極限 $H\approx 1-\varepsilon_d\xi^2$ は式~\eqref{eq:dccd_pde_form} の
PDE 形の\textbf{波数依存 $\varepsilon$} 拡張と等価で，
直流成分・6h 波長成分を $\Ord{\varepsilon_d}$ で保ちつつ
Nyquist モードのみ $4\varepsilon_d$ で減衰する．
本節の導出形 PDE は\textbf{実装デフォルトフィルタの低波数極限}であり，
式~\eqref{eq:dccd_filter_transfer} の波数選択的設計はこの極限を精緻化したものである．

% -------------------------------------------------------------------------------
\subsubsection{半離散 Fourier 安定性}
\label{sec:dccd_eigenvalue}
% -------------------------------------------------------------------------------

周期格子の対称中心 CCD 修正波数
$(D_1^{\CCD}u)_j \sim (i/\Delta x) k_{\text{eq}}(k\Delta x)\hat u$，
$(D_2^{\CCD}u)_j \sim -(1/\Delta x^2) m_{\text{eq}}(k\Delta x)\hat u$
を式~\eqref{eq:dccd_pde_form} に代入して固有値：
\begin{equation}
  \lambda(k) \;=\; -\varepsilon|c|\,\frac{m_{\text{eq}}}{\Delta x}
                 \;-\; i\,c\,\frac{k_{\text{eq}}}{\Delta x}.
  \label{eq:dccd_eigenvalue}
\end{equation}
Chu--Fan 係数 $(\beta_2,a_2,b_2)=(-1/8, 3, -9/8)$ での
修正 2 次波数 $m_{\text{eq}}(\xi)$ は
$\omega_2(\theta)^2$（式~\eqref{eq:ccd_omega_12}，$\theta=\xi$）と一致し，
小 $\xi$ 展開で $m_{\text{eq}}=\xi^2+\Ord{\xi^6}$，
Nyquist $\xi=\pi$ で有限正値 $m_{\text{eq}}(\pi)=48/5=9.6$
（$96/10$，分子 $81+48-33$・分母 $48-40+2$ の直接代入），
区間 $[0,\pi]$ で符号変化なし．従って：
\begin{equation}
  \Re\lambda(k) = -\varepsilon|c|\,\frac{m_{\text{eq}}(k\Delta x)}{\Delta x} \le 0,
  \qquad \forall k,
  \label{eq:dccd_dissipative}
\end{equation}
等号は $k=0$（直流）のみで成立．半離散系は非自明モードで無条件散逸．

\noindent\textbf{Post-filter 形での実装条件．}
本稿の DCCD は PDE 形ではなく\textbf{post-filter} 形
$H(\xi) = 1 - 4\varepsilon_d\sin^2(\xi/2)$ で実装される（式~\eqref{eq:dccd_filter_transfer}）．
フィルタ非負条件 $H(\xi)\ge 0$ から $\varepsilon_d\le 1/4$ が安定の必要条件であり，
§\ref{sec:cls_stages} で採用する $\varepsilon_d=1/4$ はこの境界値に相当する
（過剰散逸を避けつつチェッカーボードを最大限に抑える設計値）．

% -------------------------------------------------------------------------------
\subsubsection{6 つの反論への回答}
\label{sec:dccd_rebuttals}
% -------------------------------------------------------------------------------

以下の 6 反論に対し，本構成は回避または明示的にスコープアウトする：

\begin{enumerate}[label=\textbf{CA-\arabic*.}]
\item \textbf{定数 $\varepsilon$ の精度低下．}
      §\ref{sec:dccd_spectral_filter} の波数選択的フィルタが解決．
\item \textbf{境界付き領域での必要非十分性．}
      §\ref{sec:ccd_bc} の片側 Hermite 閉包が\textbf{十分条件プローブ}として機能する；
      厳密な Gustafsson--Kreiss--Sundstr\"om 安定性は
      §\ref{sec:uccd6_boundary} で開課題として明示している．
\item \textbf{境界行で $m_{\text{eq}}\ge 0$ 違反の可能性．}
      §\ref{sec:fccd_bc_options} の Option III/IV で壁近傍の行を置換．
\item \textbf{半離散 $\Re\lambda\le 0$ と完全離散 RK 安定性の齟齬．}
      TVD-RK3 の虚軸安定領域 $|\Im\lambda|\Delta t\le\sqrt{3}$ を介した CFL により，
      中程度 $\varepsilon\lesssim 0.1$ で標準 CFL に帰着．
\item \textbf{非線形・多次元・可変係数への拡張可否．}
      線形解析は\textbf{必要条件プローブ}に過ぎない；
      比較対象 WENO5（付録~\ref{app:weno5}）との相互検証で経験的裏付け．
\item \textbf{大 $\varepsilon$ の放物型時間刻み制約．}
      適応スイッチ $S(\psi)=(2\psi-1)^2$ が散逸を bulk に局在化させ，
      $\varepsilon_d\le 0.05$ の実装デフォルト設定で非拘束．
\end{enumerate}

% -------------------------------------------------------------------------------
\subsubsection{DCCD の安定化フィルタとしての役割}
\label{sec:dccd_role}
% -------------------------------------------------------------------------------

以上の導出は，実装デフォルト DCCD が\textbf{節点中心 CCD に対する外殻 post-filter} として
作用する設計であることを確立する．
旧 §~04d (\texttt{04d\_dissipative\_ccd}：DCCD スペクトル設計・スイッチ $S(\psi)$・精度影響・境界条件・保存性）は，
以下の要約として本節に吸収する：

\begin{itemize}
  \item \textbf{DCCD は移流の安定化専用}である．
        節点 CCD の散逸ゼロチャンネルへ post-filter 形式で散逸を加える設計であり，
        面位置で評価する量に対する整合性は本スキーム単独では保証されず，
        面中心化が要請される場合は FCCD（§\ref{sec:fccd_def}）が用いられる．
  \item \textbf{界面近傍は非線形スイッチで無効化}される．
        $S(\psi)=(2\psi-1)^2$ は bulk（$\psi\in\{0,1\}$）で $1$，
        界面帯（$\psi=0.5$）で $0$ を取り，
        CSF 体積力評価位置での余計な散逸を避ける．
  \item \textbf{体積保存は外殻修正}で担保．
        スペクトルフィルタ自体は \(\int u\,dx\) 不変ではないため，
        §\ref{sec:cls_advection} の CLS 体積補正（$\int\psi$ 再正規化）を要する．
        Level 1 検証（$\varepsilon_d=0.05$，$N=64$ 設定）では
        フィルタのみの体積損失 $\sim 10^{-4}\cdot T$（$T$：シミュレーション時間）が
        観測され，CLS 再正規化で機械精度まで圧縮される
        （詳細な検証値は §\ref{sec:verify_cls_advection} に提示する）．
  \item \textbf{UCCD6 との対照}．
        UCCD6（§\ref{sec:uccd6_def}）は散逸を\textbf{演算子内部}に埋め込み，
        DCCD は\textbf{外殻 post-filter} として加える．
        本稿では bulk $u, v$ 移流の既定を UCCD6 とし
        CLS 移流（$\psi$）の実装デフォルトを DCCD とする：
        Level-1 検証では UCCD6/DCCD 両形式を明示的に並列し，
        Level-2 実装既定（ch13 rising bubble）は UCCD6+AB2 を bulk に，DCCD を $\psi$ 移流に，
        Level-3 DNS は UCCD6 inside IRK を bulk に既定とする（§\ref{sec:time_int}）．
        DCCD は CLS 再初期化圧縮項の非線形フラックス散逸子
        （§\ref{sec:reinit}，UCCD6 未拡張のため）にも適用される．
\end{itemize}

\noindent
詳細な収束試験値・保存性解析・境界行設計は
§~\ref{sec:dccd_spectral_filter}（スペクトル設計）および
付録~\ref{app:ccd_analysis}（精度・楕円型解析）を参照．

% ═════════════════════════════════════════════════════════════════════════════
\subsection{DCCD パラメタの設計と適応制御}
\label{sec:dccd_params}
% ═════════════════════════════════════════════════════════════════════════════

§~\ref{sec:dccd_filter_theory} で導入した DCCD フィルタの理論的枠組みに基づき，
本節では散逸強度 $\varepsilon_{d,\max}$ の具体的な設計値と，
CLS 変数 $\psi$ に基づく適応制御アルゴリズムを与える．

\noindent\textbf{本節のスコープ．}
本節では DCCD の散逸強度パラメータ $\varepsilon_d$ と
適応制御スイッチ $S(\psi)$ の \textbf{数値設計}を扱う．
スキーム単独の数値特性（フィルタ伝達関数・安定性条件・収束性）が中心であり，
具体的な用途章割当（CLS 移流・チェッカーボード抑制・曲率帯制限など）の
機構詳細は当該章（§~\ref{sec:cls_advection} 以降）に委ねる．

表~\ref{tab:dccd_usage_map} に，本稿で DCCD が採る代表的なパラメータ設定を一覧する．

\begin{table}[htbp]
\centering
\caption{DCCD フィルタのパラメータ設計一覧（用途章割当の機構詳細は §\ref{sec:cls_advection} 以降）}
\label{tab:dccd_usage_map}
\small
\begin{tabular}{llcl}
\hline
適用シナリオ & $\varepsilon_d$ & 適応制御 & 設計根拠 \\
\hline
$C^0$ tanh プロファイル移流（界面）
  & 0.05 & あり（$S(\psi)$） & 界面プロファイル保持 \\
平滑場・チェッカーボード抑制
  & $1/4$ & なし（固定） & 安定性上限・Nyquist 完全抑制 \\
高次導関数の帯制限
  & 0.05 & あり（$S(\psi)$） & 数値振動を界面外で抑制 \\
\hline
\end{tabular}
\end{table}

% ─────────────────────────────────────────────────────────────────────────────
\subsubsection{散逸強度 \texorpdfstring{$\varepsilon_{d,\max}$}{ε\_d,max} の設計値}
\label{sec:dccd_eps_design}
% ─────────────────────────────────────────────────────────────────────────────

フィルタ伝達関数 $H(\xi;\varepsilon_d) = 1 - 4\varepsilon_d\sin^2(\xi/2)$ に対し，
低波数精度保存（波長 $\geq 6h$ で劣化 $\leq 5\%$）と
高周波抑制（ナイキスト成分 20\% 減衰）の同時要求から，
両制約が等号で一致する設計点が得られる
（スペクトル設計の詳細は §~\ref{sec:dccd_filter_theory} 参照）：
\begin{equation}
  \boxed{\varepsilon_{d,\max} = 0.05}
  \label{eq:eps_max}
\end{equation}
安定性制約 $\varepsilon_{d,\max} \leq 0.25$（式~\eqref{eq:dccd_filter_transfer}）を満足する．
より強い安定化が必要な場合は $\varepsilon_{d,\max} = 0.10$〜$0.20$ まで拡張可能だが，
低波数精度も同率で劣化する．

% ─────────────────────────────────────────────────────────────────────────────
\subsubsection{適応散逸制御則と実装アルゴリズム}
\label{sec:dccd_adaptive}
% ─────────────────────────────────────────────────────────────────────────────

\textbf{適応散逸強度．}
格子点 $i$ での散逸強度を CLS 変数 $\psi_i^n$（時刻 $t^n$）によって点別に制御する：
\begin{equation}
  \varepsilon_d^{(i)} \;=\; \varepsilon_{d,\max} \cdot S(\psi_i^n)
  \;=\; \varepsilon_{d,\max}\,(2\psi_i^n - 1)^2
  \label{eq:adaptive_eps}
\end{equation}
これにより：
\begin{itemize}
  \item $\psi_i \approx 0.5$（界面）：$\varepsilon_d^{(i)} \approx 0$，フィルタ無効，標準 CCD に退化
  \item $\psi_i \approx 0$ または $1$（バルク）：$\varepsilon_d^{(i)} \approx \varepsilon_{d,\max}$，
        最大散逸でノイズを積極的に排除
\end{itemize}

\begin{tcolorbox}[algbox, title={適応 Dissipative-CCD フィルタの実装アルゴリズム}]
\label{alg:dccd}
\textbf{入力：} スカラー場 $f$ の格子値 $\{f_i\}$，CLS 変数 $\{\psi_i^n\}$，散逸強度 $\varepsilon_{d,\max}$

\begin{enumerate}
  \item \textbf{CCD 求解：}
        標準 CCD ブロック Thomas アルゴリズム（§~\ref{sec:ccd_matrix}）を実行し，
        $\{f'_i\}$，$\{f''_i\}$ を得る．
  \item \textbf{点別散逸強度の計算：}
        各格子点 $i$ に対して
        \[
          \varepsilon_d^{(i)} = \varepsilon_{d,\max}\,(2\psi_i^n - 1)^2
        \]
        を計算する（計算コスト：$O(N)$，数値演算 1 回/点）．
  \item \textbf{1 次微分のフィルタ適用：}
        \[
          \tilde{f}'_i = f'_i + \varepsilon_d^{(i)}\,(f'_{i+1} - 2f'_i + f'_{i-1}),
          \quad i = 1, \ldots, N-2
        \]
        境界点 $i=0, N-1$：周期 BC では折り返しインデックスを使用；
        壁面・流出 BC では境界点をフィルタ適用前の CCD 値のまま保持する（§~\ref{sec:dccd_bc} 参照）．
  \item \textbf{2 次微分のフィルタ適用：}
        \[
          \tilde{f}''_i = f''_i + \varepsilon_d^{(i)}\,(f''_{i+1} - 2f''_i + f''_{i-1}),
          \quad i = 1, \ldots, N-2
        \]
  \item \textbf{後続計算への引き渡し：}
        フィルタ適用済みの $(\tilde{f}', \tilde{f}'')$ を CLS 移流，曲率計算，
        速度勾配評価に用いる．
        PPE 右辺・Corrector 発散のチェッカーボード抑制には，
        本適応制御ではなく $\varepsilon_d=1/4$ の固定フィルタ
        （§~\ref{sec:dccd_decoupling}）を用いる．
\end{enumerate}

\textbf{2 次元への拡張：}
2 次元計算では $x$ 方向と $y$ 方向に独立に 1 次元フィルタを適用する（次元分割法）：
\[
  \tilde{f}'_{x,ij} = f'_{x,ij} + \varepsilon_d^{(ij)}\,(f'_{x,i+1,j} - 2f'_{x,ij} + f'_{x,i-1,j}),
  \quad
  \tilde{f}'_{y,ij} = f'_{y,ij} + \varepsilon_d^{(ij)}\,(f'_{y,i,j+1} - 2f'_{y,ij} + f'_{y,i,j-1})
\]
$f''_{xx}, f''_{yy}, f''_{xy}$ に対しても同様に適用する．

\textbf{計算コスト：} ステップ 1（CCD，$O(N^2)$ for 2D）+ ステップ 2--4（フィルタ，$O(N^2)$）．
フィルタは CCD の後処理として追加コストが小さい．
また，$\tilde{f}'$ と $\tilde{f}''$ はフィルタ段階で\textbf{独立}に補正する——
すなわち $\tilde{f}''$ は $\tilde{f}'$ から再連立せずに計算する（計算コスト削減のための設計選択）．
\end{tcolorbox}

\noindent\textbf{適用範囲の制限．}
$\varepsilon_d^{(i)} = \varepsilon_{d,\max}(2\psi_i-1)^2$ という適応制御は
$S(\psi)$ が界面で消失する性質を持ち，
DCCD フィルタを\textbf{界面帯外に限定して適用}する自然な実装となっている．
均衡条件に依存する物理量（圧力場・毛管圧など）への適用可否は
スキーム設計の問題ではなく用途章で扱う離散整合性の問題であるため，
本節では取り扱わない．

% ═══════════════════════════════════════════════════════════════════════════════
